
\section{Evaluation}
\subsection{Experiment Setup}

Our evaluation consists of two parts. First we directly evaluate the model as
a preconditioner for solving  graph-structured linear equations using
preconditioned conjugate gradient (PCG). Second, we use the learned
preconditioning operator as a propagation operator of a shallow GCN and
evaluate whether it improves long-range propagation on long-range graph
learning tasks.

For a graph $\mathcal G$ we construct the system $K = I + \alpha L$, where $L =
    I - D^{-1/2}AD^{-1/2}$ is the normalized laplacian, and $\alpha > 0$ controls
the effective propagation range.


\paragraph{Preconditioning Training}
Both \texttt{MPSSM-PCG} and \texttt{NeuralIF} have the same training protocol.
We train the models using the objective in Eq.\ref{eq:training_loss} with 10
Gaussian probes per matrix, we also evaluate both on the test split using the
best validation checkpoint. At test time, we solve each unseen system five
times using indepenedly generated RHS. we terminate PCG (\ref{alg:pcg}) when
$\lVert r_k\rVert_2 / \lVert b\rVert_2  < \texttt{tol} = 10^{-6}$ or after 1000
iterations. Refer to the appendix for per expirement hyperparameters.

\paragraph{Long-range Propagation}

After training \texttt{MPSSM-PCG} on the dataset. We freeze the model. And
train two identical conventional GCN models, that only differ in the
propagation matrix. One time we train with $\hat A$, and the other with
$P_\theta$. For ECHO-SSSP, both models are trained using MSE (Eq.\ref{eq:mse})
and selected according to valdiation MAE (Eq.\ref{eq:mae}) at best validation
checkpoint. Refer to the appendix for per expirement hyperparameters.

Note that we do each expirement 3 times using the seeds $\{0,1,2\}$.

\subsection{Preconditioner Performence}

We first evaluate the quality of the preconditioner through its effect on the
convergence of PCG. \texttt{MPSSM-PCG} is compared agains \texttt{NeuralIF},
Jacobi and the unpreconditioned identity baseline. For each test system, we
sample $x\sim \mathcal{N}(0,I),\ b = Kx$ which provides the exact solution for
evaluating both convergence and solution accuracy. Instead of directly
constructing $M^{-1}_\theta$, we apply the preconditioner in two steps
$L_\theta y = r$ then  $L_\theta^\top z = y$. Complete residual and
solution-error statisctics are provided in the appendix.


\begin{table}[H]
    \caption{
        Preconditioning performence on ECHO-SSSP graph systems using the operator $K = I + 8L$. Mehods are
        evaluated on their PCG performence on unseen test system. Results are
        averaged over three seeds $\{0,1,2\}$.
    }
    \centering
    \small
    \label{tbl:echo_sssp_preconditioning_time}
    \begin{tabular}{c c c c c c}
        Method    & Conv. \%   $\uparrow$ & PCG Iteration $\downarrow$ & Setup Time \ref{par:setup_time} $\downarrow$ & Solve Time $\downarrow$ \ref{par:solve_time} & Total Time \ref{par:total_time}  $\downarrow$ \\
        \hline
        Identity  & $100\% \pm 0.0$       & $21.86\pm 0.00$            & $\mathbf{0.001}\pm0.000$                       & $2.225\pm0.011$                              & $2.225\pm0.011$                               \\
        Jacobi    & $100\% \pm 0.0$       & $21.50\pm 0.00$            & $0.008\pm0.001$                                & $2.240\pm0.010$                              & $2.249\pm0.010$                               \\
        NeuralIF  & $100\% \pm 0.0$       & $\mathbf{8.73} \pm 0.02$   & $0.722\pm0.001$                                & $\mathbf{1.587}\pm0.007$                     & $2.309\pm0.005$                               \\
        MPSSM-PCG & $100\% \pm 0.0$       & $8.93 \pm 0.03$            & $0.274\pm0.007$                                & $1.636\pm0.004$                              & $\mathbf{1.910}\pm0.004$                      \\
    \end{tabular}
\end{table}

The results show that both \texttt{NeuralIF} and \texttt{MPSSM-PCG} outperform
classical precondtioners in PCG iteration count, while \texttt{NeuralIF}
showing a slight advanatge over our model in that metric. However our setup
time is 2.6 times less that \texttt{NeuralIF}'s setup time, resulting in
about $17\%$ reduction in the total solve time. We see very similar results using
Peptides graphs \cite{dwivediLongRangeGraph2023} (See table \ref{tbl:peptides_struct_preconditioning_time} in appendix).

\subsection{Long-Range Propagation Performence}

Next we evaluate whether the learned preconditioner can be used as an effective long-range
propagation operator in a shallow graph neural network.

\begin{table}[H]
    \centering
    \caption{
        Effect of the propgation parameter $\alpha$ on ECHO-SSSP. The table
        shows the test MAE for and PGCN per $\alpha$ value. PGCN uses 1 layer
        for propgation. 
    }

    \label{tbl:echo_sssp_alphas}
    \begin{tabular}{c c }
        $\alpha$ & PGCN MAE   $\downarrow$ \\
        \hline
        1.0      & $0.1255 \pm 0.000571$   \\
        2.0      & $0.1174 \pm 0.000126$   \\
        4.0      & $0.1106 \pm 0.000609$   \\
        8.0      & $0.1070 \pm 0.000488$   \\
    \end{tabular}
\end{table}

In table (\ref{tbl:echo_sssp_alphas}) we perform simple abiliation to evalute
the effect of the propagation parameter $\alpha$. As established by, increasing
$\alpha$ increases the effective propgation range of $(I + \alpha L)^{-1}$. So
evaluating on a long-range dataset like ECHO-SSSP \cite{migliorCanYouHear2026} should show better
performence as $\alpha$ increases. Indeed we see that as $\alpha$ increases,
the test MAE decreases. The improvements become smaller between $\alpha = 4$
and $\alpha = 8$ which suggests, that after a certain point we get diminishing
returns.


\begin{table}[H]
    \caption{
        Depth comparison on ECHO-SSSP. The table shows the test MAE for local
        GCN and PGCN at different propagation depths. PGCN uses a learned
        preconditioning operator with $\alpha = 8.0$.
    }
    \label{tbl:echo_sssp_layers}
    \centering
    \begin{tabular}{c c c c}
        Layers & Local GCN MAE $\downarrow$ & PGCN MAE    $\downarrow$ \\
        \hline
        1      & $0.1421 \pm 0.000030$      & $0.1070 \pm 0.000488$    \\
        2      & -                          & $0.0846\pm 0.000891$     \\
        4      & $0.1112\pm 0.000242$       & $0.0791\pm 0.001215$     \\
        8      & $0.0748\pm 0.000437$       & -                        \\
    \end{tabular}
\end{table}

The second expirement aimed to test the effect of replacing $\hat A$ with
$P_\theta$. we can see from table \ref{tbl:echo_sssp_layers} that the results
are consistent: PGCN outperforms GCN with less layers. We notice that 4 layer
GCN lacks behind 1 layer of PGCN. However the gap shrinks as we add more
layers. This is partly becuase GCN applies an activiation function, and uses a
unique weight matrix for each layer. 




